% Figure IV.new --- the deontic-conflict frontier: one feature, a cliff not a
% slope. ch6-09 (register: must).
%
% Reader question: exactly which single language feature turns a free
% polynomial check into an NP-complete one?
% Claim: every rule kind already inside L_c keeps conflict detection
% polynomial with a witness (Thm. B4.1); adding ONE feature -- letting one
% obligation be disjunctive -- makes conflict-freedom NP-complete via a
% 3-SAT reduction (Thm. B4.2), and the boundary is a cliff, not a slope.
% Evidence: O(H+T log T+C log C+V*E) membership; the 3-SAT reduction (one
% disjunctive obligation per clause, two matched O/F rules per literal on
% identical scope/interval); verified on 16 seeded instances (8 SAT + 8
% UNSAT), conflict-freedom <=> satisfiability on all 16.
% Must distinguish: the identical rule vocabulary shared by both columns
% from the one added rule kind that is the entire cause of the jump.
% Grammar chosen: before/after on identical geometry (SKILL.md: "what one
% change changes") -- same vocabulary column, one row added, the verdict
% badge flips. Grammar rejected: a Venn diagram of the two fragments, which
% would suggest a fuzzy overlap where the theorem's whole point is a sharp
% boundary with no partial credit.
% Five-second test: a reader can name the one added rule kind and see the
% verdict badge flip from green to red because of it.
\begin{figure}[H]
\centering
\pdfigurehue{pdgold}%
\begin{tikzpicture}[pd figure,x=1cm,y=1cm]
  \def\lx{0.60}
  \def\rx{5.80}
  \def\dy{0.72}
  \node[pd title,anchor=west] at (\lx,4.35) {inside $L_c$};
  \node[pd title,anchor=west] at (\rx,4.35) {$L_c$ + one disjunctive obligation};
  \foreach \x in {\lx,\rx}{
    \node[pd label,anchor=west] at (\x,{4.35-1*\dy}) {Horn rules ($\to$ head, or $\bot$)};
    \node[pd label,anchor=west] at (\x,{4.35-2*\dy}) {O/F: fixed scope + interval};
    \node[pd label,anchor=west] at (\x,{4.35-3*\dy}) {exclusive interval claims};
    \node[pd label,anchor=west] at (\x,{4.35-4*\dy}) {difference constraints};}
  \node[pd breach label,anchor=west] at (\rx,{4.35-5*\dy})
    {$+O(\ell_1\vee\ell_2\vee\cdots)$: disjunctive obligation};
  \draw[pd rule,->] (\lx,{4.35-5.55*\dy}) -- (\lx,{4.35-6.35*\dy});
  \draw[pd breach rule,->] (\rx,{4.35-5.55*\dy}) -- (\rx,{4.35-6.35*\dy});
  \node[pd ready state,minimum width=4.6cm,text width=4.2cm,align=center] at ({\lx+2.2},{4.35-7.15*\dy})
    {\textbf{polynomial, witness-producing} --- three engines: Horn propagation, sweep-line overlap, Bellman-Ford};
  \node[pd breach state,minimum width=4.6cm,text width=4.2cm,align=center] at ({\rx+2.2},{4.35-7.15*\dy})
    {\textbf{NP-complete} --- selection guessed, verified by the same three engines (3-SAT reduction)};
  \node[pd note,anchor=north west,text width=9.7cm] at (\lx,{4.35-9.1*\dy})
    {verified on 16 seeded instances (8 satisfiable, 8 unsatisfiable, seed
     20260816): conflict-freedom $\Leftrightarrow$ satisfiability on all 16};
\end{tikzpicture}
\caption{Conflict checking is polynomial in the stated fragment; disjunctive obligations permit an NP-complete conflict-freedom problem. [verified, \protect\path{b4_deontic_fragment.py}]}
\label{fig:he-deontic-frontier}
\end{figure}
